% ====================================
% FORMATTING - Định dạng tài liệu
% ====================================

% Định nghĩa bảng màu đen trắng (theo yêu cầu chữ màu đen)
\definecolor{navyblue}{RGB}{0, 0, 0}          % Đen cho tiêu đề, header
\definecolor{charcoal}{RGB}{0, 0, 0}          % Đen cho chữ chính
\definecolor{lightgray}{RGB}{245, 247, 250}   % Nền xám nhạt cho code
\definecolor{bordergray}{RGB}{220, 224, 230}  % Màu viền cho bảng/code
\definecolor{linkblue}{RGB}{0, 0, 0}          % Đen cho link
\definecolor{citegreen}{RGB}{0, 0, 0}         % Đen cho trích dẫn
\definecolor{tableheader}{RGB}{224, 235, 250} % Màu nền xanh xám nhạt cho header bảng

% Thiết lập lề trang
\geometry{
    a4paper,
    left=3cm,
    right=2cm,
    top=2.5cm,
    bottom=2.5cm
}

% Giãn dòng 1.5
\onehalfspacing

% Chiều cao header tránh cảnh báo lỗi
\setlength{\headheight}{15pt}

% Thiết lập header và footer
\pagestyle{fancy}
\fancyhf{}
\fancyhead[L]{\small\color{navyblue} Hệ thống quản lý quy trình sáng tác và xuất bản Manga}
\fancyhead[R]{\small\color{navyblue} Software Engineering}
\fancyfoot[C]{\small\color{navyblue}\thepage}
\renewcommand{\headrulewidth}{0.5pt}
\renewcommand{\headrule}{{\color{navyblue}\hrule width\headwidth height\headrulewidth\hfill}}
\renewcommand{\footrulewidth}{0pt}

% Định nghĩa viết hoa cho tên chương và phụ lục
\addto\captionsvietnamese{%
  \renewcommand{\chaptername}{CHƯƠNG}%
  \renewcommand{\appendixname}{PHỤ LỤC}%
}

% Định dạng tiêu đề chương trên 1 hàng, viết hoa toàn bộ
\titleformat{\chapter}[hang]
    {\normalfont\fontsize{16pt}{20pt}\selectfont\bfseries\color{navyblue}\centering}
    {\chaptertitlename~\thechapter.~}
    {0em}
    {\MakeTextUppercase}
\titlespacing*{\chapter}{0pt}{-10pt}{30pt}

% Định dạng tiêu đề mục
\titleformat{\section}
    {\normalfont\fontsize{14pt}{18pt}\selectfont\bfseries\color{navyblue}}
    {\thesection}{1em}{}
\titlespacing*{\section}{0pt}{18pt}{10pt}

\titleformat{\subsection}
    {\normalfont\fontsize{13pt}{16.5pt}\selectfont\bfseries\color{navyblue}}
    {\thesubsection}{1em}{}
\titlespacing*{\subsection}{0pt}{14pt}{6pt}

\titleformat{\subsubsection}
    {\normalfont\fontsize{12pt}{15pt}\selectfont\bfseries\color{navyblue}}
    {\thesubsubsection}{1em}{}
\titlespacing*{\subsubsection}{0pt}{10pt}{4pt}

% Độ sâu mục lục
\setcounter{tocdepth}{2}
\setcounter{secnumdepth}{3}

% Cấu hình hiển thị tiền tố "Hình" và "Bảng" trong Danh sách hình vẽ & Danh sách bảng
\renewcommand{\cftfigpresnum}{Hình~}
\renewcommand{\cftfigaftersnum}{:}
\setlength{\cftfignumwidth}{5.0em}

\renewcommand{\cfttabpresnum}{Bảng~}
\renewcommand{\cfttabaftersnum}{:}
\setlength{\cfttabnumwidth}{5.0em}

% Thiết lập hyperref
\hypersetup{
    colorlinks=true,
    linkcolor=linkblue,
    filecolor=magenta,
    urlcolor=linkblue,
    citecolor=citegreen,
    pdftitle={Báo cáo đồ án},
    pdfauthor={Tác giả},
    bookmarks=true
}

% Thiết lập listings cho mã nguồn
\lstset{
    backgroundcolor=\color{lightgray},
    basicstyle=\ttfamily\small\color{charcoal},
    breakatwhitespace=false,
    breaklines=true,
    captionpos=b,
    commentstyle=\color{citegreen}\itshape,
    frame=single,
    rulecolor=\color{bordergray},
    numbers=left,
    numberstyle=\tiny\color{black},
    keywordstyle=\color{navyblue}\bfseries,
    stringstyle=\color{black},
    showstringspaces=false,
    tabsize=4,
    keepspaces=true
}

% Thiết lập caption
\captionsetup{
    font=small,
    labelfont=bf,
    justification=centering
}

% Định dạng đường dẫn hình ảnh
\graphicspath{{assets/images/}{assets/diagrams/}}
